% This is a modified version of Springer's LNCS template suitable for anonymized MICCAI 2025 main conference submissions.
% Original file: samplepaper.tex, a sample chapter demonstrating the LLNCS macro package for Springer Computer Science proceedings; Version 2.21 of 2022/01/12

\documentclass[runningheads]{llncs}
%
\usepackage[T1]{fontenc}
\usepackage{algorithm}
\usepackage{algorithmic}
\usepackage{amsmath,amssymb}
\usepackage{booktabs}
\usepackage{multirow}
\usepackage{array}
\usepackage{graphicx} % for \resizebox
% T1 fonts will be used to generate the final print and online PDFs,
% so please use T1 fonts in your manuscript whenever possible.
% Other font encodings may result in incorrect characters.
%
\usepackage{graphicx,verbatim}
% Used for displaying a sample figure. If possible, figure files should
% be included in EPS format.
%
% If you use the hyperref package, please uncomment the following two lines
% to display URLs in blue roman font according to Springer's eBook style:
%\usepackage{color}
%\renewcommand\UrlFont{\color{blue}\rmfamily}
%\urlstyle{rm}
%
\begin{document}
%
\title{Privacy-Preserving Federated Learning for Imbalanced Medical Applications via Distribution-Aware Aggregation}
%\titlerunning{Abbreviated paper title}
% If the paper title is too long for the running head, you can set
% an abbreviated paper title here
%
\begin{comment}  %% Removed for anonymized MICCAI submission
\author{First Author\inst{1}\orcidID{0000-1111-2222-3333} \and
Second Author\inst{2,3}\orcidID{1111-2222-3333-4444} \and
Third Author\inst{3}\orcidID{2222--3333-4444-5555}}
%
\authorrunning{F. Author et al.}
% First names are abbreviated in the running head.
% If there are more than two authors, 'et al.' is used.
%
\institute{Princeton University, Princeton NJ 08544, USA \and
Springer Heidelberg, Tiergartenstr. 17, 69121 Heidelberg, Germany
\email{lncs@springer.com}\\
\url{http://www.springer.com/gp/computer-science/lncs} \and
ABC Institute, Rupert-Karls-University Heidelberg, Heidelberg, Germany\\
\email{\{abc,lncs\}@uni-heidelberg.de}}

\end{comment}

\author{Anonymized Authors}  %% Added for anonymized MICCAI submission
\authorrunning{Anonymized Author et al.}
\institute{Anonymized Affiliations \\
    \email{email@anonymized.com}}

\maketitle              % typeset the header of the contribution
%
\begin{abstract}
Federated Learning (FL) enables collaborative model training across institutions without sharing raw data.
Performance degrades, however, when label distributions are imbalanced across sites, a common situation in medical settings where disease prevalence varies drastically.
Existing distribution-aware methods assume access to label statistics that clients are unwilling or legally unable to disclose.
We propose a model-agnostic FL framework that estimates the global label histogram once, under user-level $\varepsilon$-differential privacy, and uses the estimate to switch from standard to inverse-frequency aggregation.
Each client adds partial discrete Laplace noise to its clipped label counts and secret-shares the result via Shamir's scheme; designated decryptors reconstruct the noisy aggregate without ever observing individual contributions.
A closed-form variance decomposition, factoring clipping, client-sampling, and DP-noise components, determines the earliest round at which the estimate is reliable enough to trigger the switch.
We evaluate the framework on medical imaging and physiological signal benchmarks under realistic heterogeneous settings, demonstrating consistent improvements over standard FL aggregation with negligible communication overhead.

\keywords{Federated Learning \and Non-IID Data Distribution \and Privacy.}

\end{abstract}
%
%
%

\section{Introduction}

Federated Learning (FL)~\cite{ref_mcmahan2017} enables collaborative model training across institutions without centralising raw data, a key requirement in medical imaging where patient records cannot leave the originating hospital.
The standard aggregation rule, FedAvg~\cite{ref_mcmahan2017}, weights each client's update by its dataset size, implicitly assuming comparable label distributions across sites.
In practice, disease prevalence varies drastically: melanoma accounts for fewer than 5\% of skin lesions at general practices yet dominates specialist referrals~\cite{ref_fl_medical}, and rare arrhythmia subtypes may appear at only a handful of clinics.
The resulting label heterogeneity~\cite{ref_kairouz2021} biases the global model towards the most prevalent classes, degrading accuracy precisely on the minority conditions that are most clinically relevant.

Several methods address this imbalance by exploiting label-distribution information.
FedLA~\cite{ref_fedla} reweights aggregation using per-client label counts; FedLC~\cite{ref_fedlc} calibrates local logits by the class prior; optimisation-based approaches such as FedProx~\cite{ref_fedprox} and Scaffold~\cite{ref_scaffold} reduce client drift but do not target label skew directly.
All distribution-aware methods, however, require access to label statistics that clients are unwilling, and often legally unable, to disclose, since even label histograms can reveal sensitive population characteristics~\cite{ref_dlg,ref_label_leakage}.

This work presents a framework that bridges distribution-aware aggregation with formal privacy guarantees.
Each client adds a partial discrete Laplace noise contribution~\cite{ref_fed_heatmap} to its label counts and secret-shares the noisy result via Shamir's scheme; designated decryptors~\cite{ref_flamingo} then reconstruct the aggregate, achieving user-level $\varepsilon$-differential privacy with information-theoretic confidentiality, an approach validated at scale for federated analytics~\cite{ref_fed_heatmap}.
The histogram is estimated \emph{once}, avoiding privacy-budget composition.
Because FL already samples a random subset of clients at each round, shares accumulate naturally: we derive a closed-form variance decomposition (Eq.~\ref{eq:var}) that factors estimation error into clipping, client-sampling, and DP-noise components, allowing the server to determine the earliest round at which the estimate is reliable enough to trigger inverse-frequency aggregation.
Training therefore begins with standard FedAvg and seamlessly transitions to class-balanced weighting once sufficient clients have been observed.
The approach is model-agnostic, adds negligible communication overhead, and, unlike prior label-aware methods, never reveals individual label counts to any party.


\section{Problem statement}

We consider a $k$-class classification problem in a federated setting with $N$ clients $\mathcal{C}=\{P_1,\dots,P_N\}$ and a central server~$\mathcal{S}$.
Each client $P_i$ holds a private dataset $\mathcal{D}_i$ of size $n_i$ with label-count vector $\mathbf{y}_i=(c_{i,1},\dots,c_{i,k})\in\mathbb{Z}_{\ge0}^k$.
At each round~$r$ the server samples a subset $\mathcal{P}^r\!\subseteq\mathcal{C}$; every selected client trains locally from the current global model~$\theta^r$ to produce an update $\Delta_i^r=\theta_i^r-\theta^r$, and the server aggregates
\begin{equation}\label{eq:weighted_avg}
    \theta^{r+1} = \theta^{r} + \sum_{P_i\in\mathcal{P}^r} w_i^r\;\Delta_i^{r},
    \qquad
    w_i^r \ge 0,\;\;
    \textstyle\sum_{P_i\in\mathcal{P}^r} w_i^r = 1.
\end{equation}
Standard FedAvg~\cite{ref_mcmahan2017} sets $w_i^r = n_i / \sum_{P_j\in\mathcal{P}^r} n_j$, weighting by dataset size alone.
Under label heterogeneity this biases the model towards the most prevalent classes; distribution-aware weights can mitigate this effect, but require knowledge of the global label histogram $\mathbf{y}=\sum_i\mathbf{y}_i$, which clients are unwilling to disclose.
We formulate the following problem: \emph{estimate~$\mathbf{y}$ under $\varepsilon$-differential privacy at the user level, then use the estimate to set the weights~$w_i^r$}, without ever revealing any individual~$\mathbf{y}_i$.

\textbf{Threat model.}
Clients and a designated set of $M$~\emph{decryptors} $\mathcal{D}=\{D_1,\dots,D_M\}\subseteq\mathcal{C}$~\cite{ref_flamingo} are honest-but-curious.
The server may be malicious and may collude with up to $t-1$ decryptors, where $t=\lceil 2M/3\rceil$.

\textbf{Privacy goal.}
We provide user-level $\varepsilon$-DP for the histogram; the mechanism is detailed in Section~\ref{sec:histogram}.

\section{Method}

Our framework operates in two phases, both instances of Eq.~\eqref{eq:weighted_avg} with different weights.
During a \emph{warm-up} (rounds $1$ to $R^*$), the server uses dataset-size weights $w_i^r\!=\!n_i/\sum_j n_j$ while client coverage grows: each client generates its noisy Shamir shares once (Section~\ref{sec:histogram}) and these are buffered by the decryptors, so each additional round increases the observed fraction~$\alpha$ without additional privacy expenditure.
At round~$R^*$ the server triggers histogram reconstruction and switches to inverse-frequency weights (Section~\ref{sec:agg}).
The transition round is not a free parameter: Section~\ref{sec:trigger} derives a closed-form criterion that determines the earliest round at which the estimate is reliable enough.

\subsection{Privacy-preserving histogram estimation}\label{sec:histogram}

\textbf{Contribution bounding.}
Each client $P_i$ draws $C$ samples uniformly with replacement from~$\mathcal{D}_i$ and computes the clipped label-count vector $\tilde{\mathbf{y}}_i\!\in\!\mathbb{Z}_{\ge0}^k$ with $\|\tilde{\mathbf{y}}_i\|_1\!=\!C$, ensuring $\ell_1$-sensitivity exactly~$C$ regardless of~$n_i$.

\textbf{Client-side distributed noise.}
A P\'{o}lya random variable $X\!\sim\!\mathrm{P\acute{o}lya}(a,\beta)$ has mass $P(X\!=\!k)=\binom{a+k-1}{a-1}\beta^{k}(1-\beta)^{a}$~\cite{ref_fed_heatmap}.
By Theorem~4.1 of~\cite{ref_fed_heatmap}, if $S$~participants each independently sample $X_i,Y_i\!\overset{\mathrm{iid}}{\sim}\!\mathrm{P\acute{o}lya}(1/S,\,\beta)$ with $\beta=e^{-\varepsilon/C}$, then $Z=\sum_{i=1}^{S}(X_i-Y_i)\sim\mathrm{DLap}(C/\varepsilon)$, which achieves $\varepsilon$-DP for queries with $\ell_1$-sensitivity~$C$.

In our protocol the server broadcasts, at initialisation, the \emph{target cohort size} $|\mathcal{A}|=\alpha N$ (computed from the participation rate~$\gamma$ and the number of warm-up rounds via Eq.~\eqref{eq:trigger}).
Each client $P_i$ uses this target to calibrate its noise: for every label~$\ell$ it independently samples
\begin{equation}\label{eq:partial_noise}
    \eta_{i,\ell} = X_{i,\ell}^{+}-X_{i,\ell}^{-},\quad X_{i,\ell}^{\pm}\!\overset{\mathrm{iid}}{\sim}\!\mathrm{P\acute{o}lya}\!\bigl(\tfrac{1}{|\mathcal{A}|},\,e^{-\varepsilon/C}\bigr).
\end{equation}
When exactly $|\mathcal{A}|$~clients contribute, $\sum_{i=1}^{|\mathcal{A}|}\eta_{i,\ell}\sim\mathrm{DLap}(C/\varepsilon)$, guaranteeing $\varepsilon$-DP for the full histogram.
If fewer clients participate, the aggregated noise is smaller than $\mathrm{DLap}(C/\varepsilon)$, yielding a tighter effective privacy budget $\varepsilon_{\mathrm{eff}}>\varepsilon$; the trigger mechanism (Section~\ref{sec:trigger}) ensures reconstruction is deferred until coverage~$\alpha$ is reached.

\textbf{Shamir secret sharing.}
Each client secret-shares its \emph{noisy} count: for label~$\ell$, $P_i$ constructs a random polynomial $f_{i,\ell}$ of degree $t\!-\!1$ with $f_{i,\ell}(0)=\tilde{c}_{i,\ell}+\eta_{i,\ell}$ and sends shares $s_{i,\ell}^{(j)}=f_{i,\ell}(j)\bmod P$ to the server, which relays them to decryptor~$D_j$ ($P$ is a large prime, $t=\lceil 2M/3\rceil$).
Any subset of fewer than $t$ shares is information-theoretically independent of the secret.
Share generation is performed \emph{once} at initialisation; in subsequent rounds the server only buffers client identifiers.

\textbf{Reconstruction.}
At round~$R^*$, let $\mathcal{A}\!\subseteq\!\mathcal{C}$ be the set of clients observed so far ($|\mathcal{A}|=\alpha N$).
Each decryptor $D_j$ aggregates the shares it holds: $\bar{s}_\ell^{(j)}=\sum_{P_i\in\mathcal{A}}s_{i,\ell}^{(j)}\bmod P$.
If at least $t$ decryptors respond, the server computes Lagrange coefficients~$\lambda_j$ and reconstructs
$\hat{y}_\ell = \sum_{j=1}^{t}\lambda_j\bar{s}_\ell^{(j)}\bmod P
=\sum_{P_i\in\mathcal{A}}(\tilde{c}_{i,\ell}+\eta_{i,\ell})$,
i.e.\ the true aggregate plus the sum of client noise contributions.
If fewer than $t$ decryptors respond, estimation is deferred and FedAvg continues.

\subsection{Distribution-aware aggregation}\label{sec:agg}

Once the noisy histogram $\hat{\mathbf{y}}$ is available, the server computes class frequencies $\hat{p}_\ell=\max(\delta,\hat{y}_\ell)/\sum_{\ell'}\max(\delta,\hat{y}_{\ell'})$ ($\delta\!=\!1$ is a floor preventing division by zero) and inverse-frequency class weights
\begin{equation}\label{eq:omega}
    \omega_\ell = \frac{1/\hat{p}_\ell}{\sum_{\ell'=1}^{k}1/\hat{p}_{\ell'}}.
\end{equation}
The server broadcasts $\omega_\ell$ to all participants.
Each selected client computes a score $\phi_i = \sum_{\ell}\omega_\ell\,c_{i,\ell}$ from its own label counts and sends it together with~$\Delta_i^r$ (a single additional scalar per round).
The weights in Eq.~\eqref{eq:weighted_avg} are then set to $w_i^r = \phi_i / \sum_{P_j\in\mathcal{P}^r} \phi_j$, upweighting clients holding rare classes and downweighting those dominated by frequent ones, approximating the class-balanced loss~\cite{ref_cui2019}.
The scalar~$\phi_i$ reveals only a single weighted sum comparable to the dataset size~$n_i$ already disclosed in standard FedAvg.

\textbf{Convergence.}
Both phases use fixed weights in Eq.~\eqref{eq:weighted_avg}: Phase~1 uses dataset-size weights; Phase~2 uses inverse-frequency weights computed once at~$R^*$.
Although these weights originate from a noisy DP estimate, they are deterministic constants throughout Phase~2.
Standard convergence results for weighted FedAvg~\cite{ref_li_convergence,ref_fednova} therefore apply to each phase independently.
The warm-up serves a dual purpose: accumulating client shares for reliable histogram estimation (Section~\ref{sec:trigger}) and learning feature representations that empirically transfer to the rebalanced objective, consistent with the decoupled training paradigm~\cite{ref_kang2020}.

\subsection{Histogram uncertainty and trigger}\label{sec:trigger}

The estimate $\hat{\mathbf{y}}$ combines three error sources: partial client coverage, clipping, and DP~noise.
For a fixed label~$\ell$, let $q_{i,\ell}=c_{i,\ell}/n_i$ be its true proportion at client~$P_i$, $\bar{q}_\ell=\frac{1}{N}\sum_i q_{i,\ell}$ the unweighted mean across clients, $S_\ell^2=\mathrm{Var}(q_{i,\ell})$ the between-client variance, and $\kappa_\ell=S_\ell^2/[\bar{q}_\ell(1\!-\!\bar{q}_\ell)]\in[0,1]$ the \emph{normalised skew}.
Let $|\mathcal{A}|=\alpha N$ be the number of distinct clients observed after $R^*$ rounds.
The normalised estimation variance decomposes as
\begin{equation}\label{eq:var}
\mathrm{Var}(\hat{p}_\ell-\bar{q}_\ell)=
\frac{\bar{q}_\ell(1-\bar{q}_\ell)}{|\mathcal{A}|}
\left[
\underbrace{\frac{1-\kappa_\ell}{C}}_{\text{clipping}}
+\;
\underbrace{\kappa_\ell(1-\alpha)}_{\text{client sampling}}
\right]
+\;\sigma^2_{\mathrm{DP}},
\end{equation}
where $\sigma^2_{\mathrm{DP}}=\alpha\cdot 2\beta/(1\!-\!\beta)^2$ with $\beta=e^{-\varepsilon/C}$ is the noise variance from $|\mathcal{A}|=\alpha N$ client-side P\'{o}lya contributions (Eq.~\eqref{eq:partial_noise}); when $\alpha=1$ this equals the full $\mathrm{DLap}(C/\varepsilon)$ variance.
With $m$ clients sampled per round (with replacement across rounds), coverage reaches~$\alpha$ at round
\begin{equation}\label{eq:trigger}
R^\star=
\left\lfloor
\frac{\log(1-\alpha)}{\log\!\left(1-{m}/{N}\right)}
\right\rfloor,
\end{equation}
which the server uses to trigger reconstruction.
The bracketed term in Eq.~\eqref{eq:var} can be rewritten as $\tfrac{1}{C}+\kappa_\ell\!\bigl((1\!-\!\alpha)-\tfrac{1}{C}\bigr)$, exposing a direct trade-off between coverage~$\alpha$ (controlled via~$R^\star$) and clipping level~$C$, while $\kappa_\ell$ is dataset-dependent and not tuneable.

\section{Experiment Setup}

We evaluate the proposed aggregation method in multiple heterogeneous federated learning settings to assess its robustness across different data modalities and dimensionalities. Experiments are conducted on multiple biomedical benchmarks, including the 1D Arrhythmia dataset for time-series analysis, HAM10000 and BloodMNIST for 2D image classification, and OrganMNIST3D for volumetric 3D data. This selection enables a comprehensive evaluation of the aggregation strategy under diverse statistical properties and representation spaces, demonstrating its robustness regardless of the data modalities or chosen model.

\subsection{Type of distribution Skew}
Except for the Arrhythmia dataset, where the labels are distributed following the real partitioning assigning to each client a patient, for the other dataset we need to simulate both size skew and label skew. 

\textbf{Quantity skew (min-Dirichlet).}
Let $n_{\mathrm{tot}}$ be the total number of training samples distributed among $N$ clients.
We impose a minimum allocation $n_i \ge 0.5\,n_{\mathrm{tot}}/N$, guaranteeing at least 50\% of the uniform split per client.
The remaining samples are allocated by drawing proportions $q \sim \mathrm{Dirichlet}(\alpha_s \mathbf{1}_N)$ with $\alpha_s = 0.5$ and assigning residuals via $\mathrm{Multinomial}$.

\textbf{Label skew (Dirichlet per client).}
For each client $P_i$ we sample class proportions $\pi_i \sim \mathrm{Dirichlet}(\alpha_d \mathbf{1}_k)$ and generate label counts $\mathbf{y}_i \sim \mathrm{Multinomial}(n_i,\pi_i)$.
We denote this strategy $\mathrm{D}(\alpha_d)$: smaller $\alpha_d$ produces highly concentrated, partially missing classes; larger values approximate balanced distributions.

\subsection{Histogram Estimation}

\begin{table*}[t]
\centering
\caption{BloodMNIST distribution estimation}
\resizebox{\textwidth}{!}{%
\begin{tabular}{l c cccc cccc cc}
\toprule
\multicolumn{12}{c}{\textbf{BloodMNIST}} \\
\midrule
\multicolumn{2}{c}{\textbf{Centralized dataset}} &
\multicolumn{10}{c}{$D_{\min-s}(0.5,100,0.5)$, $\alpha=0.9$, $m=10$, $N=100$, 1000 trials} \\
\cmidrule(lr){1-2}\cmidrule(lr){3-12}
\textbf{Class} & \textbf{Real (\%)} &
\multicolumn{4}{c}{$D_l(0.1)$} &
\multicolumn{4}{c}{$D_l(0.5)$} &
\multicolumn{2}{c}{$D_l(1.0)$} \\
\cmidrule(lr){3-6}\cmidrule(lr){7-10}\cmidrule(lr){11-12}
 & &
$\bar{q}_\ell$ & Mean & Std & ExpStd &
$\bar{q}_\ell$ & Mean & Std & ExpStd &
Mean & Std \\
\midrule
basophil & 7.13 & 6.79 & 6.79 & 0.74 & 0.72 & 7.61 & 7.61 & 0.58 & 0.56 & 6.71 & 0.50 \\
eosinophil & 18.24 & 20.67 & 20.66 & 1.05 & 1.06 & 17.75 & 17.73 & 0.83 & 0.83 & 16.52 & 0.66 \\
erythroblast & 9.07 & 7.59 & 7.66 & 0.72 & 0.75 & 8.92 & 8.91 & 0.52 & 0.55 & 9.32 & 0.52 \\
granulocytes & 16.94 & 16.53 & 16.46 & 1.06 & 1.05 & 19.39 & 19.39 & 0.84 & 0.82 & 16.54 & 0.68 \\
lymphocyte & 7.10 & 6.67 & 6.67 & 0.70 & 0.71 & 8.76 & 8.77 & 0.58 & 0.61 & 8.22 & 0.53 \\
monocyte & 8.31 & 4.35 & 4.33 & 0.63 & 0.64 & 8.15 & 8.13 & 0.68 & 0.67 & 8.51 & 0.52 \\
neutrophil & 19.48 & 20.11 & 20.11 & 1.11 & 1.12 & 17.50 & 17.54 & 0.82 & 0.83 & 18.73 & 0.72 \\
platelet & 13.74 & 17.30 & 17.31 & 1.00 & 1.03 & 11.91 & 11.92 & 0.65 & 0.66 & 15.46 & 0.61 \\
\bottomrule
\end{tabular}%
}
\end{table*}



\subsection{Experimental setup}
We evaluate the proposed method under multiple experimental configurations designed to capture different sources of statistical heterogeneity. As previously discussed, for the 1D dataset the size and label skew are fixed, reflecting the natural client-wise data distribution. In contrast, for the 2D and 3D datasets, size skew is controlled using the strategy described in Section 5.1, while label skew is systematically varied through a Dirichlet partitioning with concentration parameter $\alpha_l \in {0.1, 0.5, 1}$ with a fixed client selecion scheme. This design enables a controlled comparison of aggregation behavior across diverse heterogeneity regimes.







\section{Results}


%
% ---- Bibliography ----
%
% BibTeX users should specify bibliography style 'splncs04'.
% References will then be sorted and formatted in the correct style.
%
% \bibliographystyle{splncs04}
% \bibliography{mybibliography}
%
\begin{thebibliography}{16}

\bibitem{ref_mcmahan2017}
McMahan, B., Moore, E., Ramage, D., Hampson, S., Arcas, B.A.: Communication-efficient learning of deep networks from decentralized data. In: AISTATS, pp.~1273--1282. PMLR (2017)

\bibitem{ref_kairouz2021}
Kairouz, P., McMahan, H.B., Avent, B., et~al.: Advances and open problems in federated learning. Found. Trends Mach. Learn. \textbf{14}(1--2), 1--210 (2021)

\bibitem{ref_fedprox}
Li, T., Sahu, A.K., Zaheer, M., Sanjabi, M., Talwalkar, A., Smith, V.: Federated optimization in heterogeneous networks. In: MLSys (2020)

\bibitem{ref_scaffold}
Karimireddy, S.P., Kale, S., Mohri, M., Reddi, S., Stich, S., Suresh, A.T.: Scaffold: Stochastic controlled averaging for federated learning. In: ICML, pp.~5132--5143. PMLR (2020)

\bibitem{ref_fednova}
Wang, J., Liu, Q., Liang, H., Joshi, G., Poor, H.V.: Tackling the objective inconsistency problem in heterogeneous federated optimization. In: NeurIPS, vol.~33, pp.~7611--7623 (2020)

\bibitem{ref_fedla}
Khalil, A., Wainakh, A., Zimmer, E., Parra-Arnau, J., Fernandez~Anta, A., Meuser, T., Steinmetz, R.: Label-aware aggregation for improved federated learning. In: IEEE FMEC, pp.~216--223 (2023)

\bibitem{ref_fedlc}
Zhang, J., Li, Z., Li, B., Xu, J., Wu, S., Ding, S., Wu, C.: Federated learning with label distribution skew via logits calibration. In: ICML, pp.~26311--26329. PMLR (2022)

\bibitem{ref_fl_medical}
Rieke, N., Hancox, J., Li, W., et~al.: The future of digital health with federated learning. npj Digit. Med. \textbf{3}, 119 (2020)

\bibitem{ref_fed_heatmap}
Bagdasaryan, E., Kairouz, P., Mellem, S., Gasc\'{o}n, A., Bonawitz, K., Estrin, D., Gruteser, M.: Towards sparse federated analytics: Location heatmaps under distributed differential privacy with secure aggregation. In: PoPETs, vol.~2022(4), pp.~162--182 (2022)

\bibitem{ref_dlg}
Zhu, L., Liu, Z., Han, S.: Deep leakage from gradients. In: NeurIPS, vol.~32 (2019)

\bibitem{ref_label_leakage}
Wainakh, A., Ventola, F., M\"{u}\ss ig, T., Kerstin, J., Zimmer, E., Steinmetz, R., Kersting, K.: User-level label leakage from gradients in federated learning. In: PoPETs, vol.~2022(2), pp.~227--244 (2022)

\bibitem{ref_li_convergence}
Li, X., Huang, K., Yang, W., Wang, S., Zhang, Z.: On the convergence of FedAvg on non-IID data. In: ICLR (2020)

\bibitem{ref_cui2019}
Cui, Y., Jia, M., Lin, T.-Y., Song, Y., Belongie, S.: Class-balanced loss based on effective number of samples. In: CVPR, pp.~9268--9277 (2019)

\bibitem{ref_flamingo}
Ma, Y., Woods, J., Angel, S., Polychroniadou, A., Rabin, T.: Flamingo: Multi-round single-server secure aggregation with applications to private federated learning. In: IEEE S\&P, pp.~477--496 (2023)

\bibitem{ref_kang2020}
Kang, B., Xie, S., Rohrbach, M., Yan, Z., Gordo, A., Feng, J., Kalantidis, Y.: Decoupling representation and classifier for long-tailed recognition. In: ICLR (2020)

\end{thebibliography}
\end{document}
